\chapter{Análisis del Sistema Apollo}
\label{cap:analisis_apollo}

Este capítulo describe el proceso de reverse engineering aplicado al módulo de detección de semáforos de Apollo, desde la decisión inicial de trabajar con este sistema hasta la extracción completa de las especificaciones funcionales.

\section{Selección del Sistema de Referencia}

\subsection{Contexto: El Paper de FSE 2020}

El punto de partida de este trabajo fue el paper ``A First Look at the Integration of Machine Learning Models in Complex Autonomous Driving Systems'' presentado en FSE 2020 \cite{peng2020first}. Este trabajo analiza cómo se integran los modelos de ML en Apollo, identificando:

\begin{itemize}
    \item Patrones de integración entre componentes
    \item Desafíos en el testing y validación
    \item Oportunidades para mejorar la modularidad
\end{itemize}

\textbf{Idea clave:} Los sistemas monolíticos dificultan el testing controlado de componentes individuales. Extraer subsistemas permitiría diseñar tests más específicos.

\subsection{Evaluación de Subsistemas}

Se evaluaron dos candidatos principales:

\begin{table}[H]
\centering
\begin{tabular}{|l|p{5cm}|p{5cm}|}
\hline
\textbf{Criterio} & \textbf{Detección de Semáforos} & \textbf{Detección de Obstáculos} \\
\hline
Sensores & Solo cámaras & Cámaras + LiDAR + Radar \\
\hline
Complejidad & Media & Alta \\
\hline
Reproducibilidad de tests & Alta (videos/imágenes) & Baja (requiere datos LiDAR) \\
\hline
Criticidad & Alta & Muy alta \\
\hline
Auto-contenido & Alto & Medio (depende de fusión) \\
\hline
\end{tabular}
\caption{Comparación de subsistemas candidatos}
\end{table}

\textbf{Decisión:} Se eligió el módulo de detección de semáforos por:

\begin{enumerate}
    \item \textbf{Simplicidad de testing:} Solo requiere imágenes/video como entrada
    \item \textbf{No requiere LiDAR:} Los sensores LiDAR son difíciles de simular
    \item \textbf{Interfaces claras:} Input (imagen + projection boxes) → Output (detecciones + estados)
    \item \textbf{Criticidad suficiente:} Componente de seguridad importante
\end{enumerate}

\section{Primer Enfoque: Recortar Apollo Original}

\subsection{Estrategia Inicial}

La primera aproximación fue intentar aislar el módulo TLR del código completo de Apollo:

\begin{enumerate}
    \item Clonar el repositorio completo de Apollo
    \item Identificar el módulo de percepción de semáforos
    \item Eliminar progresivamente dependencias innecesarias
    \item Mantener solo el código esencial para TLR
\end{enumerate}

\subsection{Problemas Encontrados}

Este enfoque resultó inviable por múltiples razones:

\subsubsection{Complejidad del Sistema}

\begin{itemize}
    \item \textbf{Tamaño del código:} Millones de líneas de código C++
    \item \textbf{Dependencias profundas:} El módulo TLR dependía de:
    \begin{itemize}
        \item Framework CyberRT (sistema de mensajería propio de Apollo)
        \item Librerías de utilidades internas
        \item Configuraciones específicas del ecosistema Apollo
    \end{itemize}
\end{itemize}

\subsubsection{Requisitos de Configuración}

Apollo requiere una configuración extensa y específica (ver Tabla \ref{tab:apollo_requirements}):

\begin{table}[H]
\centering
\begin{tabular}{|l|l|}
\hline
\textbf{Componente} & \textbf{Requisito} \\
\hline
CPU & 8+ cores \\
RAM & 16 GB \\
GPU & NVIDIA RTX 20xx+ \\
OS & Ubuntu 18.04/20.04/22.04 \\
CUDA & 11.8 \\
Docker & $\geq$ 19.03 \\
Espacio disco & $\sim$50 GB \\
\hline
\end{tabular}
\caption{Requisitos de hardware/software de Apollo}
\label{tab:apollo_requirements}
\end{table}

\subsubsection{Proceso de Compilación Complejo}

% TODO: Agregar diagrama del proceso de build de Apollo

La compilación de Apollo involucra:
\begin{enumerate}
    \item Descargar imagen Docker ($\sim$15 GB)
    \item Configurar container con NVIDIA runtime
    \item Compilar todo el sistema (30+ minutos)
    \item Resolver dependencias de Bazel
\end{enumerate}

\subsection{Conclusión del Primer Enfoque}

Intentar aislar el módulo TLR manteniendo el código original era:
\begin{itemize}
    \item \textbf{Técnicamente complejo:} Requería entender todo el framework CyberRT
    \item \textbf{Poco práctico:} Alta barrera de entrada para testing
    \item \textbf{Difícil de mantener:} Cualquier cambio en Apollo rompería la extracción
\end{itemize}

\textbf{Decisión:} Cambiar de estrategia a reimplementación desde cero.

\section{Segundo Enfoque: Reverse Engineering y Reimplementación}

\subsection{Nueva Estrategia}

En lugar de recortar el código existente, decidimos:

\begin{enumerate}
    \item \textbf{Analizar} el código de Apollo para entender su funcionamiento
    \item \textbf{Extraer} las especificaciones funcionales
    \item \textbf{Reimplementar} desde cero en PyTorch
    \item \textbf{Validar} equivalencia funcional
\end{enumerate}

Ventajas de este enfoque:
\begin{itemize}
    \item Control total sobre el código
    \item Dependencias mínimas
    \item Código más simple y entendible
    \item Facilita instrumentación para testing
\end{itemize}

\section{Proceso de Reverse Engineering}

\subsection{Análisis de Alto Nivel}

\subsubsection{Identificación del Pipeline}

El primer paso fue mapear el flujo completo del sistema:

\begin{figure}[H]
\centering
% TODO: Agregar diagrama del pipeline completo
\caption{Pipeline de detección de semáforos en Apollo}
\label{fig:apollo_pipeline}
\end{figure}

Etapas identificadas:
\begin{enumerate}
    \item \textbf{Preprocessing:} Obtención de ROIs y expansión de crops
    \item \textbf{Detection:} Red CNN para detectar semáforos
    \item \textbf{Assignment:} Algoritmo húngaro para asignar detecciones
    \item \textbf{Recognition:} Clasificadores especializados por orientación
    \item \textbf{Tracking:} Mantenimiento temporal de estados
\end{enumerate}

\subsubsection{Ubicación del Código}

Archivos principales del módulo TLR en Apollo:

\begin{verbatim}
modules/perception/onboard/component/
    traffic_light_region_proposal_component.cc
modules/perception/traffic_light_recognition/
    detector/detection/
        cropbox.cc              # Expansión de ROIs
        select.cc               # Algoritmo húngaro
    detector/recognition/
        classify.cc             # Clasificadores
        recognition.cc          # Lógica de reconocimiento
    tracker/
        semantic_decision.cc    # Tracking temporal
\end{verbatim}

\subsection{Extracción de Parámetros}

Se realizó un análisis sistemático del código para extraer todos los parámetros numéricos:

\subsubsection{Preprocessing}

\begin{table}[H]
\centering
\begin{tabular}{|l|c|l|}
\hline
\textbf{Parámetro} & \textbf{Valor} & \textbf{Archivo (línea)} \\
\hline
\texttt{crop\_scale} & 2.5 & cropbox.cc:28 \\
\texttt{min\_crop\_size} & 270 & cropbox.cc:29 \\
\texttt{mean\_b\_detector} & 102.98 & detection.cc:125 \\
\texttt{mean\_g\_detector} & 115.95 & detection.cc:126 \\
\texttt{mean\_r\_detector} & 122.77 & detection.cc:127 \\
\hline
\end{tabular}
\caption{Parámetros de preprocessing}
\end{table}

\subsubsection{Detection}

\begin{itemize}
    \item Input size: 270×270 (detection.cc:134)
    \item NMS threshold: 0.6 (detection.cc:401)
    \item Confidence threshold: No explícito (se acepta todo $>$ 0)
\end{itemize}

\subsubsection{Assignment (Hungarian)}

\begin{table}[H]
\centering
\begin{tabular}{|l|c|l|}
\hline
\textbf{Parámetro} & \textbf{Valor} & \textbf{Archivo (línea)} \\
\hline
\texttt{distance\_weight} & 0.7 & select.cc:71 \\
\texttt{confidence\_weight} & 0.3 & select.cc:71 \\
\texttt{gaussian\_sigma\_x} & 100.0 & select.cc:58 \\
\texttt{gaussian\_sigma\_y} & 100.0 & select.cc:58 \\
\texttt{max\_detection\_score} & 0.9 & select.cc:60 \\
\hline
\end{tabular}
\caption{Parámetros del algoritmo húngaro}
\end{table}

\subsubsection{Recognition}

\begin{table}[H]
\centering
\begin{tabular}{|l|c|c|l|}
\hline
\textbf{Clasificador} & \textbf{Input Shape} & \textbf{Threshold} & \textbf{Archivo} \\
\hline
Vertical & 96×32 & 0.5 & recognition.cc:52 \\
Horizontal & 32×96 & 0.5 & recognition.cc:52 \\
Quad & 64×64 & 0.5 & recognition.cc:52 \\
\hline
\end{tabular}
\caption{Parámetros de clasificadores}
\end{table}

Normalización para reconocimiento:
\begin{itemize}
    \item Mean BGR: [66.56, 66.58, 69.06] (classify.cc:45)
    \item Scale factor: 0.01 (classify.cc:48)
\end{itemize}

\subsubsection{Tracking}

\begin{table}[H]
\centering
\begin{tabular}{|l|c|l|}
\hline
\textbf{Parámetro} & \textbf{Valor} & \textbf{Archivo (línea)} \\
\hline
\texttt{revise\_time\_s} & 1.5 & semantic\_decision.cc:34 \\
\texttt{blink\_threshold\_s} & 0.55 & semantic\_decision.cc:35 \\
\texttt{hysteretic\_threshold} & 1 & semantic\_decision.cc:36 \\
\hline
\end{tabular}
\caption{Parámetros de tracking temporal}
\end{table}

\subsection{Extracción de Lógica de Negocio}

\subsubsection{Reglas de Tracking Temporal}

Se identificaron 5 reglas principales de seguridad:

\textbf{Regla 1: YELLOW después de RED → Mantener RED}

Código en Apollo (semantic\_decision.cc:145-149):
\begin{codesnippet}
\begin{verbatim}
if (color == YELLOW && iter->color == RED) {
    ReviseLights(lights, iter->color);
    // Mantiene RED por seguridad
}
\end{verbatim}
\end{codesnippet}

\textbf{Justificación (según comentarios en código):}
\begin{quote}
``Because of the time sequence, yellow only exists after green and before red. Any yellow after red is reset to red for the sake of safety until green displays.''
\end{quote}

\textbf{Regla 2: BLACK cuando estaba encendido → Mantener color anterior}

Apollo mantiene el último color conocido cuando detecta BLACK (apagado), asumiendo que es una detección errónea temporal.

\textbf{Regla 3: Detección de Blink (parpadeo)}

El sistema detecta el patrón BRIGHT → DARK (> 0.55s) → BRIGHT como indicador de semáforo parpadeante.

\textbf{Regla 4: Reset después de 1.5s}

Si pasan más de 1.5 segundos sin ver un semáforo, se resetea su historial y se acepta la siguiente detección sin validación.

\textbf{Regla 5: Hysteretic threshold}

Se requiere al menos 1 frame de confirmación antes de cambiar de estado.

\subsection{Análisis de la Arquitectura CNN}

\subsubsection{Detector}

El detector utiliza una arquitectura basada en Faster R-CNN con modificaciones:

\begin{itemize}
    \item \textbf{Backbone:} Red tipo VGG con capas convolucionales
    \item \textbf{RPN:} Region Proposal Network modificado (RPNProposalSSD)
    \item \textbf{ROI Pooling:} Deformable PSROIAlign
    \item \textbf{Output:} 8 valores por detección [bg, vert, quad, hori scores + 4 coords]
\end{itemize}

\subsubsection{Clasificadores}

Tres redes CNN independientes, una por orientación:

\begin{itemize}
    \item Arquitectura similar: Conv layers + FC layers
    \item Output: 4 clases [BLACK, RED, YELLOW, GREEN]
    \item Activación final: Softmax
\end{itemize}

\subsection{Conversión de Modelos}

\subsubsection{Formato Original: Caffe}

Apollo utiliza modelos en formato Caffe (.caffemodel):
\begin{itemize}
    \item Definición de arquitectura en .prototxt
    \item Pesos en formato binario .caffemodel
\end{itemize}

\subsubsection{Proceso de Conversión a PyTorch}

\begin{enumerate}
    \item \textbf{Parsear prototxt:} Extraer definición de capas
    \item \textbf{Mapear arquitectura:} Crear modelo equivalente en PyTorch
    \item \textbf{Cargar pesos:} Leer .caffemodel y convertir formato
    \item \textbf{Validar equivalencia:} Comparar outputs con misma entrada
\end{enumerate}

Herramientas utilizadas:
\begin{itemize}
    \item Scripts propios de conversión
    \item Validación numérica capa por capa
\end{itemize}

\section{Documentación Extraída}

El resultado del proceso de reverse engineering fue:

\begin{enumerate}
    \item \textbf{Tabla completa de parámetros:} Todos los valores numéricos documentados
    \item \textbf{Especificación de reglas de negocio:} Lógica de tracking y seguridad
    \item \textbf{Arquitecturas de redes:} Estructuras completas de los modelos
    \item \textbf{Flujo de datos:} Diagramas del pipeline completo
    \item \textbf{Casos especiales:} Condiciones de borde y manejo de errores
\end{enumerate}

Esta documentación sirvió como especificación funcional para la reimplementación.

\section{Hallazgos del Análisis}

\subsection{Complejidades Inesperadas}

\begin{itemize}
    \item \textbf{Triple NMS:} Se aplica NMS en tres etapas diferentes (RPN, RCNN, Global)
    \item \textbf{Ordenamiento particular:} El NMS global ordena ascendente pero procesa desde atrás
    \item \textbf{Penalizaciones en hungarian:} Detecciones fuera del crop tienen costo 0
    \item \textbf{Clipping de scores:} Los scores de detección se limitan a máximo 0.9
\end{itemize}

\subsection{Decisiones de Diseño Interesantes}

\begin{itemize}
    \item \textbf{Clasificadores especializados:} Un modelo diferente por orientación del semáforo
    \item \textbf{Enfoque basado en proyecciones:} No se busca en toda la imagen, solo en ROIs conocidos
    \item \textbf{Reglas conservadoras:} Sesgo hacia mantener RED en casos ambiguos
\end{itemize}

\subsection{Áreas Sin Documentar}

Algunos aspectos del código original carecían de documentación clara:
\begin{itemize}
    \item Justificación de valores específicos de parámetros
    \item Criterios de diseño para thresholds
    \item Casos de test utilizados durante desarrollo
\end{itemize}

Este análisis motivó la necesidad de diseñar una suite exhaustiva de tests para la reimplementación.